\documentclass[11pt]{article}
\usepackage[margin=0.78in]{geometry}
\usepackage{amsmath,amssymb,amsthm}
\usepackage[T1]{fontenc}
\usepackage{graphicx}
\usepackage[hidelinks]{hyperref}
\usepackage{microtype}

\newtheorem{theorem}{Theorem}
\newtheorem{corollary}{Corollary}
\newcommand{\PP}{\mathbb P}
\newcommand{\RF}{R^{\mathrm F}}

\title{The Bounded-Trajectory Conditioning in Conjecture 3.18 Is Null}
\author{fa\_banach\_001 / agent\_lane\_00}
\date{17 August 2026}

\begin{document}
\maketitle
\vspace{-1.1em}

\noindent\textbf{Status: candidate full negative resolution as written.}

\section{The conjecture}

Jorgensen and Pearse define a resistance network to be a connected weighted
graph with symmetric conductances $c_{zw}\geq0$ and
$c(z):=\sum_w c_{zw}<\infty$.  Its random walk has transition probabilities
$p(z,w)=c_{zw}/c(z)$.  Their Conjecture~3.18 asks for the identity
\begin{equation}\label{eq:source}
 v_x(y)=\RF(o,x)\,
 \PP_y\!\left[\tau_x<\tau_o\mid |\gamma|<\infty\right],
\end{equation}
where $|\gamma|<\infty$ means that the entire infinite trajectory is
contained in some finite subnetwork \cite[p.~22, Conj.~3.18]{JP}.

\begin{center}
\includegraphics[width=0.96\textwidth]{figures/open_problem_crop.png}
\end{center}

\section{Proof intuition}

An irreducible walk cannot be trapped forever in a finite portion of an
infinite connected graph.  From each vertex of a finite set there is a
positive-probability finite path to the outside.  Finiteness lets us choose
one common time bound and one common positive lower bound for the chance of
escape.  Repeating the estimate in time blocks makes eternal survival have
probability zero.  The event in \eqref{eq:source} is a countable union of
such null events, so ordinary conditioning on it is undefined.

\section{The null-event theorem}

Let $(X_n)_{n\geq0}$ be the network random walk and, for a vertex set $F$,
write $\tau_{F^c}:=\inf\{n\geq0:X_n\notin F\}$.

\begin{theorem}\label{thm:null}
Let $(G,c)$ be an infinite connected resistance network in the sense of
\cite{JP}.  For every starting vertex $y$ and every finite vertex set $F$,
\[
 \PP_y(X_n\in F\text{ for every }n\geq0)=0.
\]
Consequently,
\[
 \PP_y(|\gamma|<\infty)=0.
\]
\end{theorem}

\begin{proof}
If $y\notin F$, the first assertion is immediate, so suppose $y\in F$.
For every $z\in F$, connectedness and infinitude give a finite path
\[
 z=z^{(z)}_0,z^{(z)}_1,\ldots,z^{(z)}_{\ell_z}
\]
whose last vertex is outside $F$ and whose earlier vertices lie in $F$.
Every edge on this path has positive conductance and hence positive
transition probability.  Put
\[
 L:=\max_{z\in F}\ell_z,
 \qquad
 q:=\min_{z\in F}\prod_{j=0}^{\ell_z-1}
 p\bigl(z^{(z)}_j,z^{(z)}_{j+1}\bigr).
\]
The minima and maxima are over a finite set, so $L<\infty$ and $q>0$.
Starting from any $z\in F$, following the selected path shows
\[
 \PP_z(\tau_{F^c}\leq L)\geq q.
\]
Applying the Markov property at the deterministic times $L,2L,\ldots$
therefore yields, for every $m\geq1$,
\[
 \PP_y(\tau_{F^c}>mL)\leq(1-q)^m.
\]
Letting $m\to\infty$ proves that the walk cannot stay in $F$ forever.

The paper fixes an increasing exhaustion $G=\bigcup_{k\geq1}G_k$ by finite
subgraphs.  A trajectory has finite range if and only if its range is
contained in some $G_k$.  Thus
\[
 \{|\gamma|<\infty\}
 =\bigcup_{k\geq1}\{X_n\in G_k\text{ for every }n\geq0\}.
\]
Every event on the right is null by the first part, and the union is
countable.  This proves the second assertion.
\end{proof}

\begin{corollary}\label{cor:resolution}
Under the usual definition
$\PP(A\mid B)=\PP(A\cap B)/\PP(B)$ for $\PP(B)>0$, the right-hand side of
\eqref{eq:source} is undefined on every infinite resistance network.
Therefore Conjecture~3.18 is not a valid probabilistic identity as stated.
The same problem affects every conditional probability
$\PP_y[\tau_x<\tau_o\mid\gamma\subseteq G_k^{\mathrm F}]$ in the proposed
equation (3.27).
\end{corollary}

\begin{proof}
Theorem~\ref{thm:null} gives zero probability to the conditioning event in
\eqref{eq:source}.  It also gives zero probability to eternal containment in
the finite vertex set of each $G_k^{\mathrm F}$.  Both conditional
probabilities are therefore $0/0$ expressions under ordinary conditioning.
\end{proof}

\section{What a correct replacement would require}

The finite free-network law $\PP_y^{G_k^{\mathrm F}}$ is a different Markov
law, not a conditional version of $\PP_y$: deleting outward edges changes a
boundary transition from $p(z,w)=c_{zw}/c(z)$ to
\[
 p_k(z,w)=\frac{c_{zw}}{\sum_{u\in G_k}c_{zu}}.
\]
This is exactly the boundary discrepancy noted in Remark~3.19 of the source.
One may define a meaningful substitute using a specified limit of the laws
$\PP_y^{G_k^{\mathrm F}}$, a finite-time survival conditioning, a Doob
transform, or another zero-event regularization.  Such choices need not be
equivalent, and none is specified in Conjecture~3.18.  The null-event theorem
therefore resolves the statement as written while leaving open the design and
analysis of a revised probabilistic representation.

Weihrauch and Bachmann later proved that a different simple hitting-probability
representation attributed to Corollaries~3.13 and 3.15 of \cite{JP} fails in
general, and characterized its validity for transient locally finite graphs
\cite{WB}.  Their result reinforces the need to distinguish finite-network
laws from the original infinite-network law, but it does not define or repair
the null conditioning in Conjecture~3.18.

\section*{Verification, novelty, and scope}

The proof uses only connectedness, infinitude, positive transition
probability along edges, and a countable finite exhaustion; it does not
require local finiteness, transience, or recurrence.  The uniform escape
constant is valid even when vertices have infinite degree because only one
finite path per vertex of $F$ is selected.

The novelty search was bounded to the run's lightweight indexes and web/arXiv
searches for the source identifier, exact title, ``bounded trajectory'',
``finite subnetwork'', ``probabilistic representation'', and citations to the
source paper.  It found \cite{WB}, which explicitly treats a different formula,
but no later source explicitly addressing Conjecture~3.18's null event.
Novelty is therefore provisional.  Human review should verify the source
interpretation and the finite-set escape estimate, and should keep the result
scoped to ordinary conditioning; a separately specified limiting
conditioning would constitute a revised problem.

\begin{thebibliography}{9}
\bibitem{JP}
P.~E.~T. Jorgensen and E.~P.~J. Pearse,
\emph{A Hilbert space approach to effective resistance metric},
Complex Analysis and Operator Theory 4 (2010), 975--1013;
arXiv:0906.2535v3.

\bibitem{WB}
T.~Weihrauch and S.~Bachmann,
\emph{On the Probabilistic Representation of the Free Effective Resistance
of Infinite Graphs}, Journal of Theoretical Probability 36 (2023),
1956--1971; arXiv:1902.01110.
\end{thebibliography}

\end{document}
